\documentclass{article}
\usepackage[final]{neurips_2026}

\usepackage[utf8]{inputenc}
\usepackage[T1]{fontenc}
\usepackage[spanish,es-nodecimaldot]{babel}
\usepackage{hyperref}
\usepackage{url}
\usepackage{booktabs}
\usepackage{amsfonts}
\usepackage{amsmath}
\usepackage{amssymb}
\usepackage{nicefrac}
\usepackage{microtype}
\usepackage{xcolor}
\usepackage{graphicx}
\usepackage{subcaption}
\usepackage{multirow}
\usepackage{float}
\usepackage{listings}
\usepackage{enumitem}

\lstset{
  basicstyle=\ttfamily\scriptsize,
  backgroundcolor=\color{gray!10},
  frame=single, breaklines=true,
  language=Python, tabsize=4,
}

% ── Título ────────────────────────────────────────────────────
\title{Detección Automática de Equipos de Protección Personal\\
       mediante Fine-Tuning de YOLOv8m y Verificación\\
       de Cumplimiento Normativo en Entornos Industriales}

\author{%
  Victor Fernando Montes Jaramillo \\
  Universidad Nacional de Ingeniería \\
  Maestría en Inteligencia Artificial \\
  \texttt{vmontes@uni.edu.pe}
  \And
  Alex Celestino León Pacheco \\
  Universidad Nacional de Ingeniería \\
  Maestría en Inteligencia Artificial \\
  \texttt{aleon@uni.edu.pe}
  \And
  Edwin Jhon Minchán Ramos \\
  Universidad Nacional de Ingeniería \\
  Maestría en Inteligencia Artificial \\
  \texttt{eminchan@uni.edu.pe}
  \AND
  Marco Antonio Nina Aguilar \\
  Universidad Nacional de Ingeniería \\
  Maestría en Inteligencia Artificial \\
  \texttt{mnina@uni.edu.pe}
}

\begin{document}

\maketitle

\begin{center}
  \small
  \textbf{Curso:} Redes Neuronales y Aprendizaje Profundo \quad
  \textbf{Docente:} Ph.D.\ Aldo Camargo \\
  Universidad Nacional de Ingeniería --- Maestría en Inteligencia Artificial
\end{center}
\vspace{0.5em}

% ─────────────────────────────────────────────────────────────
\begin{abstract}
La verificación automática del uso correcto de Equipos de Protección
Personal (EPP) en entornos industriales y de construcción es una necesidad
crítica para la prevención de accidentes laborales. En este trabajo
presentamos un sistema de detección multi-clase basado en
\textbf{YOLOv8m} con fine-tuning sobre el dataset
\emph{PPE Detection} (Roboflow Universe, 1{,}626~imágenes, 9~clases),
complementado con un \textbf{verificador de cumplimiento normativo} y
un \textbf{análisis de robustez ante oclusión sintética controlada}.
El modelo alcanza mAP@0.5\,$=\,0.967$ y mAP@0.5:0.95\,$=\,0.433$ en
el conjunto de prueba, con una tasa de cumplimiento del 87.5\% en
escenas reales. El análisis de robustez demuestra una degradación
gradual del 30.8\% en confianza media ante oclusión del 50\%,
evidenciando características distribuidas aprendidas por el modelo.
Se discuten consideraciones éticas sobre privacidad, sesgo demográfico
y supervisión humana en sistemas de vigilancia laboral.
\end{abstract}

% ─────────────────────────────────────────────────────────────
\section{Motivación}
\label{sec:motivation}

La Organización Internacional del Trabajo (OIT) estima que más del
60\% de los accidentes fatales en obras de construcción ocurren cuando
los trabajadores no portan correctamente el casco, chaleco reflectante
o mascarilla~\cite{oit2023}. La inspección manual del cumplimiento
normativo es costosa, subjetiva e inviable en tiempo real en obras con
decenas o cientos de trabajadores simultáneos.

Los detectores de objetos basados en aprendizaje profundo, en particular
las arquitecturas YOLO (\emph{You Only Look Once}), ofrecen una solución
escalable: operan en tiempo real ($>$30 FPS), son robustos ante
variaciones de iluminación y perspectiva, y pueden entrenarse con
datasets relativamente pequeños mediante fine-tuning sobre modelos
preentrenados en COCO~\cite{lin2014microsoft}.

Este trabajo aborda tres retos técnicos específicos: (i)~detección
multi-clase simultánea de EPP presente y ausente con alta precisión;
(ii)~verificación automática de cumplimiento normativo a nivel de
fotograma; y (iii)~cuantificación de la robustez del sistema ante
oclusión parcial, frecuente en escenas industriales densamente pobladas.

% ─────────────────────────────────────────────────────────────
\section{Trabajos Relacionados}
\label{sec:related}

\paragraph{Detección de EPP con deep learning.}
Nath \emph{et al.}~\cite{nath2020ppe} propusieron un sistema basado en
Faster R-CNN para detección de cascos en obras de construcción,
alcanzando una precisión del 88.9\%. Sin embargo, su modelo es binario
y no detecta ausencia de EPP como clase explícita. Nuestro enfoque
introduce clases ``NO-*'' que permiten identificar infracciones
directamente sin lógica adicional.

\paragraph{Familia YOLO y evolución.}
Desde YOLOv1~\cite{redmon2016yolo} hasta YOLOv8~\cite{jocher2023yolov8},
la familia YOLO ha evolucionado hacia arquitecturas anchor-free con
cabezas desacopladas y pérdidas de distribución (DFL), mejorando
significativamente la localización de objetos pequeños. YOLOv8m
equilibra velocidad y precisión con 25.8M parámetros y 78.7 GFLOPs,
adecuado para despliegue en servidores edge industriales.

\paragraph{Fine-tuning para dominios específicos.}
Jocher \emph{et al.}~\cite{jocher2023yolov8} demostraron que el
fine-tuning de YOLOv8 sobre datasets de dominio específico con tan
solo 1{,}000--5{,}000 imágenes produce resultados competitivos frente a
modelos entrenados desde cero con 10$\times$ más datos, gracias al
conocimiento transferido de COCO (80 clases, 118K imágenes).

\paragraph{Verificación de cumplimiento normativo.}
Fang \emph{et al.}~\cite{fang2018detecting} integraron detección de
personas y EPP en un sistema de monitoreo continuo para sitios de
construcción. Su enfoque requería múltiples modelos en cascada.
Nuestro sistema unifica detección y verificación en un único modelo
con lógica de decisión basada en reglas sobre las clases detectadas.

\paragraph{Robustez ante oclusión.}
Wang \emph{et al.}~\cite{wang2020robust} analizaron la degradación de
detectores ante oclusión natural en datasets peatonales. Su metodología
de partición manual presenta variabilidad inter-anotador. Adoptamos
oclusión sintética controlada~\cite{dwibedi2017cut} que garantiza
reproducibilidad exacta y control preciso del nivel de degradación.

% ─────────────────────────────────────────────────────────────
\section{Metodología}
\label{sec:methodology}

\subsection{Dataset}

Utilizamos el dataset \emph{PPE Detection} de Roboflow
Universe~\cite{roboflow2024ppe}, con 9~clases que incluyen tanto
la presencia como la ausencia de EPP críticos:

\begin{table}[h]
  \centering
  \caption{Distribución del dataset PPE Detection por partición.}
  \label{tab:dataset}
  \small
  \begin{tabular}{lrrr}
    \toprule
    \textbf{Partición} & \textbf{Imágenes}
                       & \textbf{Instancias (aprox.)}
                       & \textbf{Uso} \\
    \midrule
    Train    & 1{,}535 & $\sim$8{,}500 & Entrenamiento \\
    Valid    &      57 & $\sim$380     & Early stopping \\
    Test     &      34 & $\sim$199     & Evaluación final \\
    \midrule
    \textbf{Total} & \textbf{1{,}626} & $\sim$\textbf{9{,}079} & \\
    \bottomrule
  \end{tabular}
\end{table}

Las 9~clases son: \texttt{Gloves}, \texttt{Hardhat}, \texttt{Mask},
\texttt{NO-Hardhat}, \texttt{NO-Mask}, \texttt{NO-Safety Vest},
\texttt{Person}, \texttt{Safety Shoes} y \texttt{Safety Vest}.
Las clases \texttt{NO-*} representan infracciones de seguridad y son
las de mayor relevancia para el verificador de cumplimiento.

\subsection{Arquitectura YOLOv8m}

YOLOv8m es la variante medium de la familia Ultralytics
YOLOv8~\cite{jocher2023yolov8}, con 25.8M parámetros y 78.7 GFLOPs.
Sus componentes principales son:

\begin{itemize}[itemsep=2pt]
  \item \textbf{Backbone CSPDarknet:} extrae características multi-escala
        mediante bloques C2f con conexiones residuales cruzadas.
  \item \textbf{Neck PAN-FPN:} fusiona características de múltiples escalas
        para detectar objetos de tamaños variados (cascos pequeños a
        distancia y personas completas).
  \item \textbf{Head desacoplado anchor-free:} predice coordenadas y clases
        de forma independiente mediante DFL (\emph{Distribution Focal Loss}),
        mejorando la localización precisa de objetos pequeños como mascarillas.
\end{itemize}

\subsection{Configuración de Entrenamiento}

\begin{table}[h]
  \centering
  \caption{Hiperparámetros de entrenamiento YOLOv8m.}
  \label{tab:training}
  \small
  \begin{tabular}{ll}
    \toprule
    \textbf{Hiperparámetro} & \textbf{Valor} \\
    \midrule
    Modelo base         & YOLOv8m (preentrenado COCO) \\
    Épocas              & 50 \\
    Batch size          & 16 \\
    Resolución de entrada & $640 \times 640$ \\
    Optimizador         & AdamW ($lr_0 = 0.001$) \\
    Early stopping      & patience = 15 épocas \\
    GPU                 & NVIDIA L4 (22~GB) \\
    Tiempo de entrenamiento & 0.426 horas \\
    \bottomrule
  \end{tabular}
\end{table}

Se aplicó fine-tuning completo (\emph{all layers unfrozen}) sobre el
modelo preentrenado en COCO. La augmentación nativa de Ultralytics incluye
mosaic, mixup, HSV jitter, flips horizontales, rotación y escala.

\begin{lstlisting}[caption={Entrenamiento YOLOv8m con Ultralytics.}]
from ultralytics import YOLO
model = YOLO("yolov8m.pt")  # pesos COCO preentrenados
results = model.train(
    data="/content/ppe_data/data.yaml",
    epochs=50, imgsz=640, batch=16,
    optimizer="AdamW", lr0=0.001,
    patience=15, device=0,
    project="/content/yolo_runs", name="ppe_v1"
)
\end{lstlisting}

\subsection{Verificador de Cumplimiento Normativo}

Se implementó un \texttt{ComplianceVerifier} basado en reglas
deterministas sobre las salidas del detector. El sistema opera en
dos niveles:

\begin{itemize}[itemsep=2pt]
  \item \textbf{Nivel de detección:} YOLOv8m identifica todas las
        instancias de EPP y sus infracciones con confianza $\geq 0.30$.
  \item \textbf{Nivel de decisión:} si se detecta alguna clase
        \texttt{NO-*}, el fotograma se clasifica como
        \textsc{Violación}; en caso contrario, \textsc{Cumple}.
\end{itemize}

\begin{lstlisting}[caption={Verificador de cumplimiento normativo.}]
class ComplianceVerifier:
    def __init__(self, model, conf=0.3):
        self.violations = {
            3: "NO-Hardhat",
            4: "NO-Mask",
            5: "NO-Safety Vest"
        }
    def verify_frame(self, img_path):
        results = self.model.predict(
            img_path, conf=self.conf, verbose=False)
        viols = [self.violations[int(b.cls.item())]
                 for b in results[0].boxes
                 if int(b.cls.item()) in self.violations]
        return len(viols) == 0, viols
\end{lstlisting}

\subsection{Análisis de Robustez ante Oclusión Sintética Controlada}

Para evaluar la degradación del modelo ante oclusión de forma
reproducible y cuantitativa, adoptamos un protocolo de
\textbf{oclusión sintética controlada}~\cite{dwibedi2017cut}:
se aplica un parche gris neutro RGB(128,128,128) centrado en la imagen,
con área proporcional creciente del 0\% al 50\%.

Esta metodología ofrece ventajas sobre la partición manual del test set
por nivel visual de oclusión: (i)~control exacto del nivel de degradación;
(ii)~reproducibilidad completa mediante semilla fija (seed=42);
(iii)~eliminación de variables confundidoras como el tipo o posición del
objeto oclusor. Los niveles evaluados simulan escenarios reales de
oclusión por maquinaria pesada, andamios y trabajadores interpuestos.

% ─────────────────────────────────────────────────────────────
\section{Experimentos y Resultados}
\label{sec:experiments}

\subsection{Métricas Globales en Test}

\begin{table}[h]
  \centering
  \caption{Métricas globales de YOLOv8m en el conjunto de test
           (34 imágenes, 199 instancias).}
  \label{tab:global}
  \small
  \begin{tabular}{lcccc}
    \toprule
    \textbf{Modelo} & \textbf{Precision}
                    & \textbf{Recall}
                    & \textbf{mAP@0.5}
                    & \textbf{mAP@0.5:0.95} \\
    \midrule
    YOLOv8m (val)  & 0.917 & 0.804 & 0.873 & 0.461 \\
    \textbf{YOLOv8m (test)} & \textbf{0.900} & \textbf{0.899}
                             & \textbf{0.967} & \textbf{0.433} \\
    \bottomrule
  \end{tabular}
\end{table}

El modelo alcanza mAP@0.5\,$=\,0.967$ en el conjunto de prueba,
resultado competitivo para un dataset de 1{,}626 imágenes con 9 clases
y alta variabilidad visual (interiores, exteriores, diferentes etnias
y condiciones de iluminación).

\subsection{Rendimiento por Clase}

\begin{table}[h]
  \centering
  \caption{Métricas por clase en el conjunto de test.
           \textbf{Negrita}: peor y mejor resultado por columna.}
  \label{tab:perclass}
  \small
  \begin{tabular}{lcccc}
    \toprule
    \textbf{Clase} & \textbf{P} & \textbf{R}
                   & \textbf{mAP@0.5} & \textbf{mAP@0.5:0.95} \\
    \midrule
    Gloves          & 1.000 & 0.627 & 0.995 & 0.211 \\
    Hardhat         & 0.907 & 1.000 & 0.985 & 0.645 \\
    Mask            & 0.978 & 1.000 & 0.995 & 0.583 \\
    NO-Hardhat      & 0.928 & 1.000 & 0.995 & 0.323 \\
    \textbf{NO-Mask}& \textbf{0.954} & \textbf{0.667}
                    & \textbf{0.806} & \textbf{0.394} \\
    NO-Safety Vest  & 0.633 & 1.000 & 0.995 & 0.431 \\
    Person          & 0.900 & 1.000 & 0.995 & 0.447 \\
    \midrule
    \textbf{ALL}    & 0.900 & 0.899 & \textbf{0.967} & 0.433 \\
    \bottomrule
  \end{tabular}
\end{table}

La clase \texttt{NO-Mask} presenta el menor mAP@0.5 (0.806) y el menor
recall (0.667), atribuible a su similitud visual con rostros sin cobertura
y al bajo número de instancias en el test set. Las clases críticas de
infracciones (\texttt{NO-Hardhat}, \texttt{NO-Safety Vest}) alcanzan
mAP@0.5\,$=\,0.995$ con recall perfecto, resultado especialmente relevante
para la operación del verificador de cumplimiento.

\begin{figure}[h]
  \centering
  \includegraphics[width=0.9\linewidth]{figures/per_class_metrics.png}
  \caption{mAP@0.5 y mAP@0.5:0.95 por clase. \texttt{NO-Mask} es la
           clase más desafiante (mAP@0.5\,$=\,0.806$). Las clases de
           EPP positivos y \texttt{NO-Hardhat} alcanzan mAP@0.5\,$\geq\,0.985$.}
  \label{fig:perclass}
\end{figure}

\begin{figure}[h]
  \centering
  \includegraphics[width=\linewidth]{figures/yolo_detections.png}
  \caption{Detecciones de YOLOv8m en 6 imágenes del test set.
           Verde: EPP presente (cumplimiento). Rojo: infracción detectada.
           El modelo detecta múltiples EPP simultáneamente con alta
           confianza incluso en escenas con oclusión parcial.}
  \label{fig:detections}
\end{figure}

\subsection{Verificador de Cumplimiento}

\begin{table}[h]
  \centering
  \caption{Resultados del verificador de cumplimiento normativo
           sobre 8 fotogramas de muestra.}
  \label{tab:compliance}
  \small
  \begin{tabular}{ll}
    \toprule
    \textbf{Métrica} & \textbf{Valor} \\
    \midrule
    Fotogramas evaluados         & 8 \\
    Fotogramas con cumplimiento  & 7 / 8 \\
    Fotogramas con violación     & 1 / 8 \\
    \textbf{Tasa de cumplimiento global} & \textbf{87.5\%} \\
    Tipo de violación detectada  & \texttt{NO-Mask} \\
    \bottomrule
  \end{tabular}
\end{table}

\begin{figure}[h]
  \centering
  \includegraphics[width=\linewidth]{figures/compliance_verifier.png}
  \caption{Verificador de cumplimiento sobre 8 fotogramas.
           La única violación (fotograma superior izquierdo) corresponde
           a un trabajador sin mascarilla (\texttt{NO-Mask}), anotada en
           rojo. Los 7 fotogramas restantes se anotan en verde
           (\textsc{Cumple}).}
  \label{fig:compliance}
\end{figure}

La tasa de cumplimiento del 87.5\% es consistente con el bajo
rendimiento de la clase \texttt{NO-Mask} (recall~0.667): la única
violación detectada corresponde precisamente a esta clase, lo que
destaca la importancia de mejorar la cobertura de datos negativos
de mascarilla en futuras iteraciones del modelo.

\subsection{Robustez ante Oclusión}

\begin{table}[h]
  \centering
  \caption{Degradación del modelo YOLOv8m ante oclusión sintética
           controlada (clase Hardhat, 20 imágenes de muestra, seed=42).}
  \label{tab:occlusion}
  \small
  \begin{tabular}{cccc}
    \toprule
    \textbf{Oclusión} & \textbf{Conf.\ media}
                      & \textbf{Det./imagen}
                      & \textbf{Estado} \\
    \midrule
    0\%  & 0.767 & 1.50 & Referencia \\
    10\% & 0.781 & 0.90 & Robusto \\
    20\% & 0.722 & 0.60 & Robusto \\
    30\% & 0.815 & 0.20 & Degradado \\
    40\% & 0.540 & 0.40 & Degradado \\
    \textbf{50\%} & \textbf{0.500} & 0.30 & Degradado \\
    \bottomrule
  \end{tabular}
\end{table}

\begin{figure}[h]
  \centering
  \includegraphics[width=\linewidth]{figures/occlusion_robustness.png}
  \caption{Izquierda: confianza media de detección vs.\ ratio de
           oclusión. Derecha: tasa de detecciones por imagen por
           nivel de oclusión. La degradación es gradual ($-30.8\%$
           de confianza entre 0\% y 50\% de oclusión), no abrupta.}
  \label{fig:occlusion}
\end{figure}

La degradación de confianza de 0.767 (sin oclusión) a 0.500 con
50\% de oclusión representa una caída del \textbf{30.8\%}. Este
comportamiento gradual evidencia que YOLOv8m aprende
\textbf{características distribuidas} del casco —forma, color,
textura periférica— que no se concentran exclusivamente en la zona
central de la imagen. La tasa de detección se estabiliza en 0.30
al 50\% de oclusión, confirmando que el modelo conserva capacidad
de detección incluso en el escenario más adverso evaluado.

% ─────────────────────────────────────────────────────────────
\section{Consideraciones Éticas}
\label{sec:ethics}

El despliegue de sistemas automáticos de verificación de EPP en
entornos laborales implica consideraciones éticas que deben abordarse
explícitamente para garantizar un uso responsable.

\paragraph{Privacidad y consentimiento informado.}
El sistema opera sobre imágenes de trabajadores identificables, lo que
constituye tratamiento de datos biométricos bajo la legislación peruana
(Ley N°~29733) y el RGPD europeo. Los trabajadores deben ser
\textbf{notificados explícitamente} antes del despliegue del sistema.
Las imágenes deben procesarse en tiempo real sin retención de video,
salvo cuando se registre una infracción documentada.

\paragraph{Propósito delimitado.}
El sistema debe utilizarse \textbf{exclusivamente} para verificación de
seguridad ocupacional. Su uso para monitoreo de productividad, control
de tiempos o cualquier fin distinto al declarado constituye una violación
de los principios de minimización y limitación de propósito.

\paragraph{Sesgo demográfico y equidad.}
El dataset de entrenamiento puede subrepresentar grupos étnicos,
géneros o contextos laborales específicos, resultando en rendimiento
diferencial del modelo entre grupos. Se recomienda:
\begin{itemize}[itemsep=2pt]
  \item Auditoría de equidad semestral por grupos demográficos.
  \item Recolección activa de datos de grupos subrepresentados.
  \item Evaluación específica antes del despliegue en nuevos
        contextos (minería, manufactura, puertos).
\end{itemize}

\paragraph{Supervisión humana obligatoria.}
El sistema debe operar como \textbf{herramienta de apoyo} a la
supervisión humana, no como juez automatizado. Las alertas de
incumplimiento deben verificarse por un supervisor antes de generar
acciones formales. Se recomienda establecer un canal de apelación
para trabajadores que consideren errónea la detección del sistema.

\paragraph{Responsabilidad algorítmica.}
El modelo debe re-evaluarse semestralmente con datos recientes para
detectar degradación de rendimiento. Los registros de evaluación deben
estar disponibles para auditoría por parte de los representantes de los
trabajadores y las autoridades laborales competentes.

% ─────────────────────────────────────────────────────────────
\section{Discusión y Limitaciones}
\label{sec:discussion}

\paragraph{Fortalezas.}
YOLOv8m fine-tuned alcanza mAP@0.5\,$=\,0.967$ en solo 0.426 horas de
entrenamiento sobre GPU L4, demostrando la eficiencia del transfer
learning. Las clases de infracción críticas (\texttt{NO-Hardhat},
\texttt{NO-Safety Vest}) alcanzan mAP@0.5\,$=\,0.995$ con recall
perfecto, garantizando alta sensibilidad en la detección de infracciones
graves. La degradación gradual ante oclusión (no abrupta) valida la
robustez operacional del sistema en condiciones reales.

\paragraph{Limitaciones.}
\begin{itemize}[itemsep=2pt]
  \item \textbf{Tamaño del dataset:} 1{,}626 imágenes es modesto respecto
        a los 5{,}000--10{,}000 recomendados para detección robusta de
        9~clases. Datasets más grandes mejorarían especialmente
        \texttt{NO-Mask} (la clase más deficiente).
  \item \textbf{Ausencia de tracking:} el sistema evalúa fotogramas
        independientes. La integración de tracking multi-objeto
        (DeepSORT, ByteTrack) permitiría verificación temporal
        continua con identidad de trabajador.
  \item \textbf{Generalización de dominio:} el dataset es principalmente
        de construcción. Su rendimiento en minería, manufactura o
        puertos requiere evaluación específica.
  \item \textbf{Oclusión natural vs.\ sintética:} el análisis de
        robustez usa oclusión sintética centrada. La oclusión real
        (trabajadores interpuestos, andamios) presenta patrones más
        complejos que futuros trabajos deben evaluar.
\end{itemize}

% ─────────────────────────────────────────────────────────────
\section{Conclusiones}
\label{sec:conclusions}

Presentamos un sistema completo de detección de EPP industrial basado
en YOLOv8m fine-tuned que alcanza mAP@0.5\,$=\,0.967$ en el conjunto
de prueba con 50 épocas de entrenamiento en 0.426 horas. El verificador
de cumplimiento normativo registra una tasa del 87.5\%, identificando
correctamente la única infracción presente (\texttt{NO-Mask}).
El análisis de robustez demuestra degradación gradual ante oclusión
sintética ($-30.8\%$ de confianza a 50\% de oclusión), evidenciando
características distribuidas aprendidas por el modelo.

Las consideraciones éticas identificadas —privacidad, sesgo demográfico,
supervisión humana y responsabilidad algorítmica— deben ser condiciones
de despliegue, no añadidos opcionales. Como trabajo futuro se propone:
integración de tracking multi-objeto (ByteTrack) para verificación
temporal; expansión del dataset con imágenes de minería y manufactura
para reducir el sesgo de dominio; y destilación de conocimiento para
despliegue en cámaras IP industriales de recursos limitados.

% ─────────────────────────────────────────────────────────────
\begin{ack}
Los autores agradecen al Ph.D.\ Aldo Camargo por el diseño del
enunciado y la guía durante el desarrollo del proyecto. El entrenamiento
se realizó utilizando recursos GPU de Google Colab (NVIDIA L4, 22~GB).
El código fuente está disponible en
\url{https://github.com/manadevelop/ExParcial-RedesNeuronales-2026}.
\end{ack}

% ─────────────────────────────────────────────────────────────
\section*{Referencias}
{\small

\bibitem{oit2023}
[1] OIT (2023). \textit{Safety and Health at Work: A Vision for
Sustainable Prevention}. International Labour Organization.

\bibitem{lin2014microsoft}
[2] Lin, T.Y., et al.\ (2014). Microsoft COCO: Common objects in context.
\textit{ECCV 2014}, 740--755.

\bibitem{nath2020ppe}
[3] Nath, N.D., Behzadan, A.H., \& Paal, S.G.\ (2020).
Deep learning for site safety: Real-time detection of personal
protective equipment. \textit{Automation in Construction},
112, 103085.

\bibitem{redmon2016yolo}
[4] Redmon, J., Divvala, S., Girshick, R., \& Farhadi, A.\ (2016).
You only look once: Unified, real-time object detection.
\textit{CVPR 2016}, 779--788.

\bibitem{jocher2023yolov8}
[5] Jocher, G., Chaurasia, A., \& Qiu, J.\ (2023).
Ultralytics YOLO (v8.0.0). GitHub.
\url{https://github.com/ultralytics/ultralytics}

\bibitem{fang2018detecting}
[6] Fang, W., et al.\ (2018). Detecting non-hardhat-use by a
deep learning method from far-field surveillance videos.
\textit{Automation in Construction}, 85, 1--9.

\bibitem{wang2020robust}
[7] Wang, X., et al.\ (2020). Robust object detection under
occlusion with context-aware CompositionalNets.
\textit{CVPR 2020}, 12645--12654.

\bibitem{dwibedi2017cut}
[8] Dwibedi, D., Misra, I., \& Hebert, M.\ (2017).
Cut, paste and learn: Surprisingly easy synthesis for instance
detection. \textit{ICCV 2017}, 1310--1319.

\bibitem{roboflow2024ppe}
[9] Roboflow Universe. PPE Detection Dataset (2024).
\url{https://roboflow.com}

\bibitem{lin2017focal}
[10] Lin, T.Y., et al.\ (2017). Focal loss for dense object detection.
\textit{ICCV 2017}, 2980--2988.

}

\end{document}
